
\begin{equation}\tag{1.12.1}
  \lambda_t(E).\sigma_{-t}(E)=1.
\end{equation}

Cette relation elle-même montre l'additivité de l'application \(\sigma_t\), qui peut donc être prolongée en un homomorphisme de \(K^{\bullet}(S)\) dans \(1+K^{\bullet}(S)[[t]]^+\)~:

\textbf{Remarque 1.13.} On peut prendre pour base de \(K^{\bullet}(X)\) sur \(K^{\bullet}(S)\) n'importe quelle suite de puissances de \(\xi\)~; \(\xi^k,\ldots,\xi^{k+r}\), car \(\xi\) est inversible, et une telle suite est donc composée d'éléments linéairement indépendants sur \(K^{\bullet}(S)\); par ailleurs, les autres puissances de \(\xi\) peuvent s'exprimer explicitement comme combinaisons linéaires à coefficients dans \(K^{\bullet}(S)\) de celles-ci, d'après l'égalité (1.9.1), et l'inversibilité de \(\xi\).

Posons \(x=1-\xi^{-1}\) (\(x\) s'interprète comme «classe d'une section hyperplane»). On peut alors prendre pour base de \(K^{\bullet}(X)\) sur \(K^{\bullet}(S)\) les éléments \(1,x,\ldots,x^r\). Pour obtenir l'équation de dépendance de \(1,\ldots,x^{r+1}\) on considère la suite exacte :
\[
  0\longrightarrow \underline{F}\longrightarrow f^*(\underline{E})\longrightarrow \underline{O}_X(1)\longrightarrow0,
\]
et la suite exacte obtenue par passage aux duals :
\[
  0\longrightarrow \underline{O}_X(-1)\longrightarrow f^*(\check{\underline{E}})\longrightarrow \check{\underline{F}}\longrightarrow0.
\]
On remarque alors que \(\gamma^{r+1}(\check F-r)=0\), d'où l'on déduit :
\[
  \gamma^{r+1}(\check F-r)=\gamma^{r+1}(\check E-r-1+x)=0,
\]
soit :
\begin{equation}\tag{1.13.1}
  \sum_{k=0}^{r+1}\gamma^{r+1-k}(\check E-r-1).x^k=0.
\end{equation}

On obtient un résultat analogue en introduisant \(x'=\xi-1\).

Notons enfin que si \(\underline{E}=(\underline{O}_S)^{r+1}\), donc \(X=\mathbf P_S^r\), on a \(E=r+1\), et l'égalité (1.13.1) se réduit à :
